%% kapa: exam=1 round=29 year=2018 period=2 form=A subject=appraisal_law subject_ko=감정평가관계법규 seq=2 number=2
%% source: Q-Net 감정평가사 기출문제 (cst00309) artlSeq=5202808 file=extracted/1008/5202808/2018년 제29회 감정평가사 1차시험 기출문제/2018년 제29회 감정평가사 1차시험 기출문제/2교시 A형.pdf
%% conversion: kordoc 4.24.51 -> tools/gen1.py (auto); review=proofread-sonnet; flags=none
\begin{question}{2}
국토의 계획 및 이용에 관한 법령상 도시ㆍ군관리계획에 관한 설명으로 옳지 않은 것은?
\begin{choices}
  \item 도시ㆍ군관리계획 결정의 효력은 지형도면을 고시한 날의 다음날부터 발생한다.
  \item 주민은 기반시설의 설치ㆍ정비 또는 개량에 관한 사항에 대하여 도시ㆍ군관리계획의 입안을 제안할 수 있다.
  \item 도시ㆍ군관리계획의 입안 시 주민의 의견을 청취하여야 하는 경우 그에 필요한 사항은 대통령령이 정하는 기준에 따라 해당 지방자치단체의 조례로 정한다.
  \item 국가계획과 관련되어 국토교통부장관이 입안한 도시ㆍ군관리계획은 국토교통부장관이 결정한다.
  \item 도시ㆍ군관리계획을 조속히 입안하여야 할 필요가 있다고 인정되면 광역도시계획이나 도시ㆍ군기본계획을 수립할 때에 도시ㆍ군관리계획을 함께 입안할 수 있다.
\end{choices}
\end{question}
\answer{1}
